\documentclass[12pt]{article}
\usepackage[utf8]{inputenc}
\usepackage{amsmath, amssymb, mathtools}
\usepackage{geometry}
\geometry{margin=1in}

\title{Quiz 2 -- Solutions\\ \vspace{8pt} \Large ECON 308 -- International Economic Relations}
\author{Instructor: Muhebullah Karimzada}
\date{October 02,  2025}


\begin{document}
\maketitle

\section*{Answers to Multiple-Choice Questions}

\noindent 1. \textbf{B) Capital owners.} \\
When $P_M/P_F$ rises, labor moves toward $M$. The value marginal product of $K$ rises both because of the higher $P_M$ and because $L_M$ increases. Landowners lose; workers are ambiguous.

\bigskip

\noindent 2. \textbf{C) More labor moves into manufactures.} \\
Equilibrium requires $w=P_M \cdot MPL_M(L_M)=P_F \cdot MPL_F(L_F)$. A higher $P_M$ shifts up the $M$-schedule, so equality holds only with larger $L_M$ (and smaller $L_F$).

\bigskip

\noindent 3. \textbf{C) Because real wage rises in terms of one good and falls in terms of the other.} \\
Typically $w$ rises less than $P_M$, so $w/P_M$ falls, but $w/P_F$ rises when $P_F$ is unchanged.

\bigskip

\noindent 4. \textbf{C) Specific factors in the export sector gain, specific factors in the import sector lose, labor is ambiguous.} \\
Classic Ch.~4 distributional result: the specific factor in the sector whose relative price rises gains; the other specific factor loses; labor’s real income is mixed.

\bigskip

\noindent 5. \textbf{B) Home’s opportunity cost of manufactures is lower than Foreign’s.} \\
Home exports $M$ under free trade only if $OC_M^{Home}<OC_M^{Foreign}$ (equivalently, $P_M/P_F$ under autarky is lower at Home).

\bigskip
\hrule
\bigskip

\section*{Solution to the Problem}


Two goods, $M$ and $F$. Labor $L$ mobile across sectors; $K$ specific to $M$; $T$ specific to $F$. \\
Production:
\[
Q_M = K^{1/3} L_M^{2/3}, 
\qquad 
Q_F = T^{1/2} L_F^{1/2}, 
\qquad 
L_M+L_F=L.
\]
Endowments: $K=64$, $T=100$, $L=90$. \\
Prices: initially $P_M=2$, $P_F=1$. Later $P_M$ rises to $3$ (with $P_F=1$).

\bigskip

\noindent \textbf{1. Marginal products of labor.}

\[
MPL_M \equiv \frac{\partial Q_M}{\partial L_M} 
= \frac{2}{3}\, K^{1/3} L_M^{-1/3}, 
\qquad
MPL_F \equiv \frac{\partial Q_F}{\partial L_F} 
= \frac{1}{2}\, T^{1/2} L_F^{-1/2}.
\]

\noindent Both are decreasing in own labor, so moving labor into a sector lowers its marginal product.

\bigskip

\noindent \textbf{2. Initial equilibrium $(L_M,L_F,w)$ with $P_M=2$, $P_F=1$.}

\emph{Wage equalization (labor mobility):}
\[
w \;=\; P_M \cdot MPL_M(L_M) \;=\; P_F \cdot MPL_F(L_F), 
\qquad L_F=90-L_M.
\]
Plug constants $K^{1/3}=4$, $T^{1/2}=10$:
\[
2 \cdot \frac{2}{3}\cdot 4 \cdot L_M^{-1/3}
\;=\; 1 \cdot \frac{1}{2}\cdot 10 \cdot (90-L_M)^{-1/2}.
\]
Simplify:
\[
\frac{16}{3} \, L_M^{-1/3} \;=\; 5 \, (90-L_M)^{-1/2}.
\]
This equation is monotone in $L_M$ and has a unique interior solution. Solving numerically:
\[
\boxed{\,L_M \approx 74.446,\quad L_F \approx 15.554\,}.
\]
The common wage is (either side of the equality):
\[
w \;=\; P_M \cdot MPL_M(L_M)
= \frac{16}{3}\, L_M^{-1/3}
\;\Rightarrow\;
\boxed{\,w \approx 1.268\,}.
\]

\bigskip

\noindent \textbf{3. Outputs and specific-factor returns $(Q_M,Q_F,r_K,r_T)$.}

\emph{Outputs:}
\[
Q_M = K^{1/3} L_M^{2/3} = 4\,L_M^{2/3} \;\Rightarrow\; \boxed{\,Q_M \approx 70.787\,},
\]
\[
Q_F = T^{1/2} L_F^{1/2} = 10\,\sqrt{L_F} \;\Rightarrow\; \boxed{\,Q_F \approx 39.438\,}.
\]

\emph{Specific-factor returns:}
\[
r_K = P_M \cdot \frac{\partial Q_M}{\partial K}
= 2 \cdot \Big(\tfrac{1}{3} K^{-2/3} L_M^{2/3}\Big).
\]
Since $K^{2/3}=16$, $K^{-2/3}=1/16$, so
\[
r_K = 2 \cdot \frac{1}{3}\cdot \frac{1}{16} \, L_M^{2/3}
= \frac{1}{24}\, L_M^{2/3}
\;\Rightarrow\; \boxed{\,r_K \approx 0.737\,}.
\]
For land:
\[
r_T = P_F \cdot \frac{\partial Q_F}{\partial T}
= 1 \cdot \Big( \tfrac{1}{2} T^{-1/2} L_F^{1/2} \Big)
= \frac{1}{20}\, \sqrt{L_F}
\;\Rightarrow\; \boxed{\,r_T \approx 0.197\,}.
\]

\noindent With these endowments and prices, most labor goes to $M$. The high $L_M$ raises $M$'s marginal productivity of $K$, so $r_K>r_T$.

\bigskip
\newpage
\noindent \textbf{4. Effects of a rise in $P_M$ from $2$ to $3$ (qualitative, then numeric check).}

\emph{Comparative statics (signs):}
\[
w = P_M \cdot MPL_M(L_M) = P_F \cdot MPL_F(L_F).
\]
When $P_M$ increases:
\begin{itemize}
\item The $M$-side wage schedule shifts up $\Rightarrow$ \textbf{$L_M$ rises}, $L_F$ falls.
\item Because $L_M$ rises, $MPL_M$ falls, so \textbf{$w$ rises less than proportionally} to $P_M$.
\item \textbf{$r_K$ rises more than proportionally}: it scales with $P_M$ and with $L_M^{2/3}$, which also rises.
\item \textbf{$r_T$ falls}: $L_F$ falls, lowering $MPL_T$; $P_F$ is unchanged.
\end{itemize}

\emph{Numerical check (optional but informative):} new equilibrium solves
\[
3 \cdot \frac{2}{3}\cdot 4 \, L_M^{-1/3} \;=\; 5 \,(90-L_M)^{-1/2}
\;\;\Rightarrow\;\;
8\, L_M^{-1/3} = 5 \,(90-L_M)^{-1/2}.
\]
Solution:
\[
\boxed{\,L_M' \approx 82.592,\quad L_F' \approx 7.408\,},\qquad
\boxed{\,w' = 8\,L_M'^{-1/3} \approx 1.837\,}.
\]
Specific-factor returns:
\[
r_K' = P_M' \cdot \frac{\partial Q_M}{\partial K}
= 3 \cdot \frac{1}{48} L_M'^{2/3}
= \frac{1}{16} L_M'^{2/3}
\;\Rightarrow\; \boxed{\,r_K' \approx 1.185\,},
\]
\[
r_T' = \frac{1}{20}\sqrt{L_F'}
\;\Rightarrow\; \boxed{\,r_T' \approx 0.136\,}.
\]
Percent changes:
\[
\Delta P_M = +50\%,\quad 
\Delta w \approx +44.9\% \ (<50\%),\quad
\Delta r_K \approx +60.8\%,\quad
\Delta r_T \approx -30.9\%.
\]
\textit{Interpretation:} Capital owners gain the most; landowners lose; workers’ nominal wage rises but less than $P_M$.

\bigskip

\noindent \textbf{5. Labor’s real wage and the “ambiguity.”}

Initial real wages:
\[
\frac{w}{P_M} \approx \frac{1.268}{2} = 0.634,
\qquad
\frac{w}{P_F} \approx \frac{1.268}{1} = 1.268.
\]
After the shock:
\[
\frac{w'}{P_M'} \approx \frac{1.837}{3} = 0.612 \quad \text{(falls)},
\qquad
\frac{w'}{P_F} \approx \frac{1.837}{1} = 1.837 \quad \text{(rises)}.
\]
\textit{Conclusion:} Labor’s purchasing power falls in terms of $M$ but rises in terms of $F$. Without a specified utility function (or consumption basket), the net welfare effect for workers is \emph{ambiguous}, consistent with the core prediction of the specific-factors model in Krugman--Obstfeld--Melitz (Ch.~4).

\end{document}
